\chapter{Analyse et Étude de Fonctions}

\section{Sandbox}
\subsection{subsection}

\begin{df}[Plop]
	Mot d'origine inconnu qui envoi du pâté ... enfin on pense : $$e^{i\pi}+1=0$$
\end{df}

\begin{thm}[Théorème des valeurs intermédiaires]
	Soit $f$ une fonction continue sur $\intff{a}{b}$ et un réel $k$ tel que $f(a)\lt k \lt f(b)$.

	Il existe au moins un réel $\alpha\in\intff{a}{b}$ tel $f(\alpha)=k$.
\end{thm}


\begin{prop}[Interéssant]
	Les chat retombrent toujours sur leurs pattes
\end{prop}

\begin{demo}\lipsum[1]\end{demo}

\begin{meth}[Calculer une intégrale]
	Faut trouver une primitive de $f$ sur l'intervalle puis faire $F(b)-F(a)$.
	$$\dint_a^b f(x)\, \dx = \Big\lbrack F(x)\Big\rbrack_a^b = F(b) - F(a)$$
\end{meth}


\begin{ex}
	Soit $f(x)=x^2$ définie sur $\intff{0}{2}$.

	On a donc : $F(x)=\frac{x^3}{3}+K$ avec $K\in\R$.

	$$
		\def\arraystretch{2}
		\begin{array}{rcl}
			\dint_0^1 f(x)\,\dx & = & \dint_0^1 x^2\,\dx \\ &=& \bra{\dfrac{x^3}{3}+K}_0^1 \\ &=& \pa{\dfrac{1^3}{3}+K}-\pa{\dfrac{0^3}{3}+K} =\dfrac13
		\end{array}
	$$
\end{ex}

\begin{rem}
	L'aire définie entre \ldots

	\bi
	\item \ldots la courbe \Cf représentative de la fonction $x^2$,
	\item l'axe des abscisses,
	\item l'axe des ordonnées,
	\item la droite d'éqiation $x=1$\ldots
	\ei

	\ldots a pour valeur $\dfrac13$.

	\begin{tikzpicture}[x=2cm,y=2cm,>=Stealth]
		\draw[dotted,xstep=0.5,ystep=0.5] (-0.3,-0.3) grid (1.3,1.3);
		\draw[->] (-0.3,0) -> (1.3,0);
		\draw[->] (0,-0.3) -> (0,1.3);
		\foreach \pt in { (0,0) , (0,1) , (1,1) , (1,0) } {\filldraw \pt circle (1pt);}
		\draw[thick,color=black,smooth,samples=100,domain=0:1] plot(\x,{\x*\x});
	\end{tikzpicture}
\end{rem}

\begin{exo}
	Soit la fonction $f(x)=\cos(x)$ définie sur $\intff{0}{2\pi}$.

	\begin{enumerate}
		\item Déterminer une primitive de $f$.
		\item Calculer :
			$$\dint_0^{\frac{\pi}{2}} f(x)\,\dx$$
	\end{enumerate}
\end{exo}


\section{Ensembles de nombres et Généralités}

\begin{df}
	On a : \np{12345} \epsdice{1}

	Dans ce chapitre, nous travaillerons principalement dans l'ensemble des nombres réels $\R$, mais rappelons les inclusions usuelles :
	\[
		\N \subset \Z \subset \D \subset \Q \subset \R \subset \C \quad \text{avec } \np{12345.678} \in \R \text{ et } \eur \in \text{Monnaie}
	\]
\end{df}

\begin{definition}
	Soit $f$ une fonction définie sur un intervalle $I = \intfo{0}{+\infty}$ et un plan doté d'un repère orthonormé $\vOIJ$.

	L'ensemble des points $M$ de coordonnées $\coordl{x}{f(x)}$ est la courbe représentative de la fonction $f$ et elle est notée $\Cf$.
	\begin{center}
		\begin{tikzpicture}[>=Stealth]
			\draw[dotted] (-7.5,-1.5) grid (7.5,1.5);
			\draw[->,thick,black!60] (0,-1.5) -- (0,1.5);
			\draw[->,thick,black!60] (-7.5,0) -- (7.5,0);
			\foreach \x in {-7,...,7} {\filldraw (\x,0) circle (1pt);}
			\foreach \y in {-1,0,1} {\filldraw (0,\y) circle (1pt);}
			\filldraw[red] (0,0) circle (2pt);
			\draw[line width=0.5pt,color=black,smooth,dashed,samples=1000,domain=-7.5:7.5] plot(\x,{sin(\x*\x*180/pi)});
		\end{tikzpicture}
	\end{center}

\end{definition}

\begin{propriete}
	Pour tout $x \in \intoo{-1}{1}$, si $f(x) = \abs{x} + 1$, alors la valeur absolue vérifie l'inégalité triangulaire :

	$$\abs{x+y}\le\abs{x}+\abs{\frac35}\le a\ge b\gt c\lt d$$

	\[
		\norm{\vec{u} + \vec{v}} \le \norm{\vec{u}} + \norm{\vec{v}}
	\]

	On note également la relation logique $\eqv$ (ou $\iff$) et l'implication $\Rarr$.
\end{propriete}

\begin{demo}
	Soient $a, b \in \intff{a}{b}$. On procède par double implication :

	\bi
	\item Sens direct : supposons $a \le b$.
	\item Sens indirect : supposons l'inverse.
	\ei

	\lipsum[2]
\end{demo}

\begin{demo}
	\lipsum[1]
\end{demo}

\section{Méthodes et Exemples}

\begin{theoreme}[Calcul d'une intégrale]
	Pour calculer l'intégrale d'une fonction continue sur $\intff{a}{b}$, on utilise la notation :
	\[
		I = \dint_a^b f(x) \dx\dy\dz
	\]
\end{theoreme}

\begin{exemple}
	Considérons la fonction $f(x) = \bfrac{x+1}{2}$.
\end{exemple}

\begin{remarque}
	On note $\pdx{f}$ et $\pdy{f}$ les dérivées partielles de $f$.
	$$
		\begin{array}{rrcl}
			      & x   & = & 3 \\
			\Rarr & x^2 & = & 9
		\end{array}
	$$
\end{remarque}

\begin{exercice}
	Soit $f$ une fonction. Résoudre dans $\intof{-\pi}{\pi}$ l'équation $f(x) = 0$.
\end{exercice}

\chapter{Géométrie Vectorielle et Algorithmique}

\section{Repères et Représentation Matrix}

Soit un repère orthonormé $\vOIJ$ (ou $\OIJ$).

\begin{theoreme}
	Considérons les points $A$ et $B$ de coordonnées respectives $\coord{x_A}{y_A}$ et $\coord{x_B}{y_B}$.
	Le vecteur $\vec{AB}$ a pour coordonnées $\coord{x_B - x_A}{y_B - y_A}$.
\end{theoreme}

\begin{thm}
	Deux droites $d_1$ et $d_2$ sont parallèles, noté $d_1 \parallel d_2$, si et seulement si le déterminant des vecteurs directeurs est nul :
	\[
		\determinant{a & b \\ c & d} = 0 \iff ad - bc = 0
	\]
\end{thm}

Nous pouvons également manipuler des matrices et des systèmes :
\[
	M = \bmat{1 & 2 \\ 3 & 4} \quad \text{ou sous forme réduite} \quad \pmat{a & b \\ c & d}
\]

$$\pfrac{1}{2}+\bfrac{1}{2}=\brac{1;2;3}=1,234=\np{1444441.234}$$

\section{Algorithmique en Python}

Pour illustrer les compétences d'algorithmique du programme, voici une fonction Python simple avec gestion des accents :

\begin{lstlisting}[language=Python]
def vecteur_colineaire(u, v):
    # Verification de la colinearite de deux vecteurs u et v
    det = u[0] * v[1] - u[1] * v[0]
    if det == 0:
        print("Les vecteurs sont colinéaires !")
        return True
    else:
        print("Les vecteurs ne sont pas colineaires.")
        return False
\end{lstlisting}

\section{Questions et Exercices Récapitulatifs}

\bmc{2} % Début du mode multicolonnes
\be % Début de l'énumération raccourcie
\item Calculer la somme $\np{1500.6789} \eur + \np{5,4343} \eur$.
\item Calculer la norme $\norm{\vec{u}}$.
\item Calculer la norme $\norm{\vec{u}}$.
\item Calculer la norme $\norm{\vec{u}}$.
\item Calculer la norme $\norm{\vec{u}}$.
\item Vérifier si $\coordl{2}{3} \in \Cf$.
\ee
\emc

On a \Cf \Cg \Ch prout

\section{symboles}
\noindent\faBook\,\faWarning\,\faRadiation\,\faYoutube
\,\faIcon{500px}
\,\faAccessibleIcon
\,\faAccusoft
\,\faAcquisitionsIncorporated
\,\faAd
\,\faAddressBook
\,\faAddressBook[regular]
\,\faAddressCard
\,\faAddressCard[regular]
\,\faAdjust
\,\faAdn
\,\faAdversal
\,\faAffiliatetheme
\,\faAirbnb
\,\faAirFreshener
\,\faAlgolia
\,\faAlignCenter
\,\faAlignJustify
\,\faAlignLeft
\,\faAlignRight
\,\faAlipay\\
\,\faAllergies
\,\faAmazon
\,\faAmazonPay
\,\faAmbulance
\,\faAmericanSignLanguageInterpreting
\,\faAmilia
\,\faAnchor
\,\faAndroid
\,\faAngellist
\,\faAngleDoubleDown
\,\faAngleDoubleLeft
\,\faAngleDoubleRight
\,\faAngleDoubleUp
\,\faAngleDown
\,\faAngleLeft
\,\faAngleRight
\,\faAngleUp
\,\faAngry
\,\faAngry[regular]
\,\faAngrycreative\\
\,\faAngular
\,\faAnkh
\,\faApper
\,\faApple
\,\faApple*
\,\faApplePay
\,\faAppStore
\,\faAppStoreIos
\,\faArchive
\,\faArchway
\,\faArrowAltCircleDown
\,\faArrowAltCircleDown[regular]
\,\faArrowAltCircleLeft
\,\faArrowAltCircleLeft[regular]
\,\faArrowAltCircleRight
\,\faArrowAltCircleRight[regular]
\,\faArrowAltCircleUp
\,\faArrowAltCircleUp[regular]
\,\faArrowCircleDown
\,\faArrowCircleLeft
\,\faArrowCircleRight
\,\faArrowCircleUp
\,\faArrowDown
\,\faArrowLeft
\,\faArrowRight
\,\faArrows*
\,\faArrowsAltH
\,\faArrowsAltV
\,\faArrowUp
\,\faArtstation\\
\,\faAssistiveListeningSystems
\,\faAsterisk
\,\faAsymmetrik
\,\faAt
\,\faAtlas
\,\faAtlassian
\,\faAtom
\,\faAudible
\,\faAudioDescription
\,\faAutoprefixer
\,\faAvianex
\,\faAviato
\,\faAward
\,\faAws
\,\faBaby
\,\faBabyCarriage
\,\faBackspace
\,\faBackward
\,\faBacon
\,\faBacteria
\,\faBacterium
\,\faBahai
\,\faBalanceScale
\,\faBalanceScaleLeft
\,\faBalanceScaleRight
\,\faBan
\,\faBandAid
\,\faBandcamp
\,\faBarcode
\,\faBars\\
\,\faBaseballBall
\,\faBasketballBall
\,\faBath
\,\faBatteryEmpty
\,\faBatteryFull
\,\faBatteryHalf
\,\faBatteryQuarter
\,\faBatteryThreeQuarters
\,\faBattleNet
\,\faBed
\,\faBeer
\,\faBehance
\,\faBehanceSquare
\,\faBell
\,\faBell[regular]
\,\faBellSlash
\,\faBellSlash[regular]
\,\faBezierCurve
\,\faBible
\,\faBicycle
\,\faBiking
\,\faBimobject
\,\faBinoculars
\,\faBiohazard
\,\faBirthdayCake
\,\faBitbucket
\,\faBitcoin
\,\faBity
\,\faBlackberry
\,\faBlackTie\\
\,\faBlender
\,\faBlenderPhone
\,\faBlind
\,\faBlog
\,\faBlogger
\,\faBloggerB
\,\faBluetooth
\,\faBluetoothB
\,\faBold
\,\faBolt
\,\faBomb
\,\faBone
\,\faBong
\,\faBook
\,\faBookDead
\,\faBookmark
\,\faBookmark[regular]
\,\faBookMedical
\,\faBookOpen
\,\faBookReader
\,\faBootstrap
\,\faBorderAll
\,\faBorderNone
\,\faBorderStyle
\,\faBowlingBall
\,\faBox
\,\faBoxes
\,\faBoxOpen
\,\faBoxTissue
\,\faBraille\\
\,\faBrain
\,\faBreadSlice
\,\faBriefcase
\,\faBiking
\,\faBimobject
\,\faBinoculars
\,\faBiohazard
\,\faBirthdayCake
\,\faBitbucket
\,\faBitcoin
\,\faBity
\,\faBlackberry
\,\faBlackTie
\,\faBlender
\,\faBlenderPhone
\,\faBlind
\,\faBlog
\,\faBlogger
\,\faBloggerB
\,\faBluetooth
\,\faBluetoothB
\,\faBold
\,\faBolt
\,\faBomb
\,\faBone
\,\faBong
\,\faBook
\,\faBookDead
\,\faBookmark
\,\faBookmark[regular]\\
\,\faBookMedical
\,\faBookOpen
\,\faBookReader
\,\faBootstrap
\,\faBorderAll
\,\faBorderNone
\,\faBorderStyle
\,\faBowlingBall
\,\faBox
\,\faBoxes
\,\faBoxOpen
\,\faBoxTissue
\,\faBraille
\,\faBrain
\,\faBreadSlice
\,\faBriefcase

\ms

Pour avoir la doc : \lstinline{texdoc fontawesome5}
